% ============================================================
\section{Discussion}

\subsection{Final Output vs. Internal Exposure}

Internal leakage was consistently higher than final-output leakage across all four systems. Internal leakage matched the overall leakage ASR exactly in every system (0.958 to 1.000), while final-output leakage ranged from 0.083 to 0.458. The resulting output-only miss rate ranged from 0.542 to 0.875, meaning that between roughly half and seven-eighths of successful leaks would have been invisible to an evaluation that inspected only the final response. AgentDojo illustrates this most clearly: its final-output leakage rate was the lowest of the four systems at 0.083, yet its internal leakage rate was 0.958, nearly identical to the other three systems. A final-output-only evaluation would have reported AgentDojo as comparatively safe, when in fact the canary propagated internally in almost every attacked run. Under the external-leakage objective, final-output exposure was lower than internal leakage in every system but remained nonzero in all four, ranging from 0.167 to 0.375. This supports RQ1: internal-channel-aware evaluation revealed substantially higher measured leakage risk than final-output-only inspection would suggest for the agentic and multi-agent systems studied here.

The qualitative pattern observed in our experiments---substantially greater leakage in monitored internal channels than in final user-visible outputs---is consistent with the internal-channel findings reported by AgentLeak~\cite{el2026agentleak}. This provides external support for the importance of internal-channel-aware evaluation and for the argument that final-output-only evaluation can miss security-relevant information propagation. The absolute leakage rates should not be directly compared across the two studies, because they differ in models, tasks, attack construction, and instrumentation; this comparison supports the importance of internal-channel measurement only.

\subsection{Performance Degradation}

RQ2 asks whether leakage-oriented attacks can be extended to degrade task
utility or increase operational burden. Automated search did not identify an
attack that reduced task success in the reported degradation experiment. A
separate diagnostic evaluation indicated that manually specified degradation
mechanisms can affect task success and operational burden, but the automated
search did not reproduce that behavior.

\subsection{Variation Across Frameworks}

RQ3 asks how leakage and final-output exposure vary across frameworks, with
degradation treated as a secondary effect. Leakage varied little: all four
systems reached a leakage ASR between 0.958 and 1.000. Final-output exposure
varied across systems, from 0.083 in the general leakage setting on AgentDojo
to 0.375 under the external-leakage objective on LangGraph, while remaining
nonzero in every system under the external-leakage objective. Given the small
number of attacked runs per configuration, these cross-system differences
should be read as descriptive patterns rather than definitive rankings. The
automated degradation search produced no measurable task-success drop in its
reported experiment.

% ============================================================
\section{Limitations}

This work has several limitations, spanning the scale of the evaluation, the
design of the search, and what was not evaluated.

\textbf{Evaluation scale.} Results come from a small-scale exploratory
evaluation rather than a definitive benchmark. Each small-scale configuration
used four tasks and six attack variants, which may introduce variance at the
family level. The evaluation used a single local model, \texttt{llama3.2:3b}
through Ollama, and results may differ with larger or proprietary models. The
medium-scale degradation result was conducted only for LangGraph, using a
generic degradation setup rather than the targeted family sweep, and should
be interpreted as a pilot rather than a general finding.

\textbf{Search design and budget.} The AutoResearch search explored only 6 of
the 1,928 known attack configurations per small-scale run (12 for the
medium-scale pilot). Of those 6 iterations, only 4 had a prior score to react
to, since crossover requires at least two previously scored variants. The
reported results therefore reflect a small, mostly cold-start search budget
rather than a converged search. Search quality more broadly depends on the
defined attack space and scoring function, and a different space or reward
weighting could surface different attack variants.

\textbf{Detection method.} Some semantic leakage may be difficult to detect
using exact or partial canary matching alone, since a model can paraphrase or
partially reconstruct a secret without reproducing it verbatim.

\textbf{Scope of comparison.} This work does not include a random or manually
designed attack baseline, a cross-framework transfer evaluation, or a defense
comparison.

\textbf{CrewAI external-leakage baseline.} The lower clean success rate
observed for CrewAI in the external-leakage setting should be interpreted
cautiously due to the small number of tasks and possible grading variance.

\textbf{Deployment realism.} Canary secrets are synthetic markers used only
for measurement; they are not real credentials, private data, or protected
production secrets. A detected canary exposure therefore measures
propagation through the experimental workflow, not a real-world breach. As
defined in Section~7.1, canary presence in the original attacker-controlled
input is not counted as leakage.

\textbf{Resource and time constraints.} The evaluation was constrained by
limited compute, time, and search budget. Each small configuration used only
four tasks and six attack variants, and the AutoResearch search was mostly
cold-start rather than converged. Future work should evaluate larger task sets,
larger search budgets, multiple seeds and models, stronger baselines, and
broader framework or runtime coverage. Where repeated runs are available,
future work should also report confidence intervals or other uncertainty
estimates rather than relying on single descriptive proportions.

% ============================================================
\section{Ethical Considerations}

All experiments are conducted in controlled environments using synthetic data. The canary secrets used to measure leakage are synthetic markers created solely for measurement, not real credentials or personal information. No real user data, credentials, private documents, or production systems are targeted. The goal of this work is to improve how agentic and multi-agent systems are evaluated for information leakage under adversarial conditions, with performance degradation considered as a secondary extension.

\paragraph{AI Usage Disclosure.}
AI-based tools were used as supporting assistants during the development of this project. Specifically, they were used for language refinement, wording improvements, and spelling and grammar checks. AI assistance was also used during the software development process to support code implementation, debugging, and the configuration and execution of experimental setups. In addition, the live experiments used \texttt{llama3.2:3b} through Ollama as the evaluated model inside the agentic workflows.

All AI-assisted outputs were reviewed and validated by the project authors. The research problem formulation, methodological design, experimental decisions, analysis, interpretation of results, and conclusions remain the responsibility and original work of the authors.
% ============================================================
\section{Conclusion}

This project developed an AutoResearch-inspired red-teaming framework for agentic and multi-agent LLM systems, in which a heuristic search procedure generates, executes, and scores attack variants against a target system. We evaluated the framework on four systems, LangGraph, AgentDojo, AutoGen, and CrewAI, with information leakage as the primary objective and final-output exposure as a leakage-specific comparison point. Performance degradation was evaluated as a secondary extension rather than as a co-equal objective. Leakage was highly effective across all four systems, with attack success rates between 0.958 and 1.000, and internal-channel leakage consistently exceeded final-output leakage, so an evaluation that inspected only the final output would have substantially underestimated measured exposure. External leakage ASR was lower than internal leakage in every system but remained nonzero in all four. Automated search did not identify an attack that reduced task success in the reported degradation experiment, although the separate diagnostic evaluation indicated that degradation-inducing attack mechanisms exist in the tested setting. The automated search did not reproduce the behavior observed in that diagnostic experiment. Together, the results show that channel-aware evaluation can reveal propagation that is invisible in final-output-only measurements, while the limited degradation search highlights the need for larger and better-controlled evaluations of secondary degradation effects.
